\documentclass[aspectratio=169,11pt]{beamer}

\usepackage[T1]{fontenc}
\usepackage{lmodern}
\usepackage{booktabs}
\usepackage{algorithm2e}
\input{../common/hkmacros.tex}

% ------------------------------------------------------------------ theme
\usetheme{default}
\setbeamercolor{structure}{fg=hkblue}
\setbeamercolor{frametitle}{fg=hkblue}
\setbeamercolor{block title}{bg=hkblue!14,fg=hkblue!85!black}
\setbeamercolor{block body}{bg=hkblue!5}
\setbeamercolor{block title alerted}{bg=hkred!14,fg=hkred!90!black}
\setbeamercolor{block body alerted}{bg=hkred!5}
\setbeamercolor{alerted text}{fg=hkred}
\setbeamerfont{frametitle}{size=\Large,series=\bfseries}
\setbeamerfont{block title}{size=\normalsize,series=\bfseries}
\setbeamertemplate{navigation symbols}{}
\setbeamertemplate{blocks}[rounded][shadow=false]
\setbeamertemplate{itemize items}[circle]

\setbeamertemplate{frametitle}{%
  \vskip1.2ex
  \hspace*{0.2cm}\usebeamercolor[fg]{frametitle}\usebeamerfont{frametitle}\insertframetitle\par
  \vskip-0.4ex
  \hspace*{0.2cm}{\color{hkblue!55}\rule{\dimexpr\paperwidth-0.6cm\relax}{0.6pt}}%
}

\newcommand{\curspk}{A}
\setbeamertemplate{footline}{%
  \leavevmode\hbox to\paperwidth{%
    \hspace{0.3cm}{\footnotesize\color{hkgrey}Hopcroft--Karp \textperiodcentered\ CS-G526}%
    \hfill{\footnotesize\color{hkgrey}Presenter \curspk\quad \insertframenumber\,/\,\inserttotalframenumber}\hspace{0.3cm}}%
  \vskip1.2ex}

\title{Bipartite Matching in $O(m\sqrt n)$ Time}
\author{Presenter A \and Presenter B}
\date{CS-G526 \textperiodcentered\ November 2026}

\DontPrintSemicolon
\SetAlgoLined
\SetKwComment{Comment}{// }{}

% layered picture for phase 2 of the running example  (M = a1, b2, c4)
\newcommand{\layerA}{%
  \node[font=\footnotesize\bfseries,text=hkblue] at (0,1.75) {layer 0};
  \node[font=\footnotesize\bfseries,text=hkblue] at (2.2,1.75) {layer 1};
  \node[font=\footnotesize\bfseries,text=hkblue] at (4.4,1.75) {layer 2};
  \node[font=\footnotesize\bfseries,text=hkblue] at (6.6,1.75) {layer 3};
  \node[boy]  (Ld) at (0,0.8)  {$d$};
  \node[boy]  (Le) at (0,-0.8) {$e$};
  \node[girl] (L1) at (2.2,0.8)  {$1$};
  \node[girl] (L2) at (2.2,-0.8) {$2$};
  \node[boy]  (La) at (4.4,0.8)  {$a$};
  \node[boy]  (Lb) at (4.4,-0.8) {$b$};
  \node[girl] (L3) at (6.6,0.8)  {$3$};
  \node[girl] (L5) at (6.6,-0.8) {$5$};
}

\begin{document}

% ======================================================================
\begin{frame}[plain]
\begin{tikzpicture}[remember picture,overlay]
  \fill[hkblue!6] (current page.south west) rectangle (current page.north east);
  \fill[hkblue] (current page.north west) rectangle ([yshift=-3mm]current page.north east);
\end{tikzpicture}
\vspace{0.5cm}
\begin{columns}[c,onlytextwidth]
\column{0.62\textwidth}
  {\color{hkblue}\Huge\bfseries Bipartite Matching\\[1mm] in $O(m\sqrt n)$ Time\par}
  \vspace{4mm}
  {\Large\color{black!70} The Hopcroft--Karp Algorithm\par}
  \vspace{10mm}
  {\large Presenter A \ \& \ Presenter B\par}
  \vspace{1mm}
  {\normalsize CS-G526 \textperiodcentered\ November 2026\par}
  \vspace{9mm}
  {\footnotesize\color{hkgrey}J.\,E.\ Hopcroft and R.\,M.\ Karp,
   \emph{An $n^{5/2}$ algorithm for maximum matchings in bipartite graphs},
   SIAM J.\ Comput.\ 2(4):225--231, 1973.\par}
\column{0.36\textwidth}
\centering
\begin{tikzpicture}[scale=0.8,transform shape]
  \exnodes \exedges
  \draw[mat] (xa)--(y3); \draw[mat] (xb)--(y5); \draw[mat] (xc)--(y4);
  \draw[mat] (xd)--(y2); \draw[mat] (xe)--(y1);
\end{tikzpicture}
\end{columns}
\end{frame}

% ======================================================================
\begin{frame}{The story in four questions}
\Large
\begin{enumerate}\setlength\itemsep{12pt}
  \item[\textcolor{hkblue}{\textbf{A1}}] \textbf{What} is the problem?\\
  {\normalsize\color{hkgrey}Jobs and workers. Why does the obvious greedy idea fail?}
  \item[\textcolor{hkblue}{\textbf{A2}}] \textbf{How} do we repair a bad assignment?\\
  {\normalsize\color{hkgrey}Chain reactions of swaps (augmenting paths). When can we stop?}
  \item[\textcolor{hkorange}{\textbf{B1}}] \textbf{How} do we speed it up?\\
  {\normalsize\color{hkgrey}Do many swaps at once: ripples (BFS) and a careful search (DFS).}
  \item[\textcolor{hkorange}{\textbf{B2}}] \textbf{How fast} is it really?\\
  {\normalsize\color{hkgrey}Why only about $\sqrt n$ rounds are ever needed.}
\end{enumerate}
\end{frame}

% ======================================================================
\begin{frame}{Workers and jobs}
\begin{columns}[c]
\column{0.45\textwidth}
\centering
\begin{tikzpicture}[scale=1.1]
  \exnodes \exedges
  \node[above=2mm of xa, font=\small\bfseries, text=hkblue] {workers};
  \node[above=2mm of y3, font=\small\bfseries, text=hkorange] {jobs};
\end{tikzpicture}
\column{0.53\textwidth}
\Large
Each line means: \emph{this worker can do this job.}\\[4mm]
\normalsize
\begin{itemize}\setlength\itemsep{8pt}
  \item A worker does \textbf{at most one} job.
  \item A job needs \textbf{at most one} worker.
\end{itemize}
\vspace{3mm}
\begin{alertblock}{Question}
What is the largest number of jobs we can fill at the same time?
\end{alertblock}
\end{columns}
\end{frame}

% ======================================================================
\begin{frame}{Same thing, in the language of graphs}
\begin{columns}[T]
\column{0.5\textwidth}
\begin{block}{Bipartite graph}
Two groups (workers $X$, jobs $Y$). Edges only go \emph{between} the groups.
\end{block}
\begin{block}{Matching}
A set of edges where \textbf{no two share a vertex}: each worker and job used at most once.
\end{block}
\begin{block}{Free vertex}
A worker or job that currently has no partner.
\end{block}
\column{0.46\textwidth}
\begin{alertblock}{The problem}
Find a matching with the \textbf{most edges} (a \emph{maximum} matching).
\end{alertblock}
\vspace{3mm}
Where it shows up: assigning tasks, scheduling, pairing students with projects, and as a building block in bigger algorithms.
\vspace{4mm}

\textbf{Before 1973:} $O(mn)$ time.\\
\textbf{Hopcroft--Karp:} $O\big((m+n)\sqrt n\big)$.
\end{columns}
\end{frame}

% ======================================================================
\begin{frame}{First idea: be greedy}
\begin{columns}[c]
\column{0.44\textwidth}
\centering
\begin{tikzpicture}[scale=1.05]
  \exnodes \exedges
  \draw[mat] (xa)--(y1);
  \only<2->{\draw[mat] (xb)--(y2);}
  \only<3->{\draw[mat] (xc)--(y4);}
  \only<4->{\freering{xd}\freering{xe}\freering{y3}\freering{y5}}
\end{tikzpicture}
\column{0.54\textwidth}
Take an edge whenever both ends are free.
\begin{itemize}\setlength\itemsep{6pt}
  \item<1-> $a\!-\!1$
  \item<2-> $b\!-\!2$
  \item<3-> $c\!-\!4$
  \item<4-> Now $d$ wants $\{1,2\}$ and $e$ wants $\{1\}$: all taken. \alert{Stuck.}
\end{itemize}
\vspace{2mm}
\onslide<5->{
Three pairs, yet five are possible: $e1,\,d2,\,b5,\,c4,\,a3$.
\begin{alertblock}{The real question}
Can we \emph{repair} the assignment without throwing it away?
\end{alertblock}}
\end{columns}
\end{frame}

% ======================================================================
\begin{frame}{The repair: a chain reaction}
\begin{columns}[c]
\column{0.44\textwidth}
\centering
\begin{tikzpicture}[scale=1.05]
  \exnodes \exedges
  \draw[mat] (xb)--(y2); \draw[mat] (xc)--(y4);
  \only<1-4>{\draw[mat] (xa)--(y1);}
  \only<5->{\draw[mat] (xa)--(y3); \draw[mat] (xd)--(y1);}
  \only<1>{\freering{xd}\freering{y3}}
  \only<2-4>{\draw[aug] (xd)--(y1); \freering{xd}\freering{y3}}
  \only<3-4>{\draw[aug] (xa)--(y1);}
  \only<4>{\draw[aug] (xa)--(y3);}
\end{tikzpicture}
\column{0.54\textwidth}
\begin{enumerate}\setlength\itemsep{7pt}
  \item<1-> Worker $d$ is unassigned. Job $3$ is unfilled.
  \item<2-> $d$ asks for job $1$.
  \item<3-> Job $1$ belongs to $a$. Ask $a$ to give it up\ldots
  \item<4-> \ldots $a$ can take job $3$, which is free!
  \item<5-> Everyone is happy: $d\!-\!1$ and $a\!-\!3$. \textbf{One more pair!}
\end{enumerate}
\vspace{2mm}
\onslide<5->{\begin{alertblock}{}
We re-routed, not restarted. Matching size: $3\to4$.
\end{alertblock}}
\end{columns}
\end{frame}

% ======================================================================
\begin{frame}{That chain has a name: an augmenting path}
\begin{block}{Augmenting path}
Start at a \textbf{free} worker. Alternate: edge \emph{not} in the matching, edge \emph{in} the matching, edge \emph{not} in\ldots{} End at a \textbf{free} job.
\end{block}
\vspace{2mm}
\centering
\begin{tikzpicture}[scale=1.0]
  \node[boy] (a) at (0,0) {$d$};
  \node[girl] (b) at (2.4,0) {$1$};
  \node[boy] (c) at (4.8,0) {$a$};
  \node[girl] (d) at (7.2,0) {$3$};
  \draw[aug] (a)--node[above,font=\footnotesize]{not in $M$}(b);
  \draw[aug,line width=3.4pt] (b)--node[above=1pt,font=\footnotesize]{in $M$}(c);
  \draw[aug] (c)--node[above,font=\footnotesize]{not in $M$}(d);
  \freering{a}\freering{d}
  \node[font=\footnotesize,text=hkgreen] at (0,-0.9) {free};
  \node[font=\footnotesize,text=hkgreen] at (7.2,-0.9) {free};
\end{tikzpicture}

\vspace{4mm}
\raggedright
\textbf{Flip it:} in-matching edges leave, the others enter.
\begin{itemize}\setlength\itemsep{4pt}
  \item One more edge enters than leaves $\Rightarrow$ the matching grows by \textbf{exactly 1}.
  \item Middle vertices just swap partner; both ends get a partner. Nobody is double-booked.
\end{itemize}
\end{frame}

% ======================================================================
\begin{frame}{When can we stop?}
\begin{alertblock}{Berge's theorem}
A matching is maximum \;$\Longleftrightarrow$\; there is \textbf{no augmenting path}.
\end{alertblock}
\vspace{2mm}
{\small\textbf{Why?} Growing is clear. For the other direction, compare $M$ with a \emph{better} matching $N$ by drawing both on top of each other:}

\begin{center}
\begin{tikzpicture}[scale=1.0,v/.style={circle,fill=black!70,inner sep=1.5pt}]
  \node[v](p1) at (0,0){}; \node[v](p2) at (1.2,0){}; \node[v](p3) at (2.4,0){}; \node[v](p4) at (3.6,0){};
  \draw[Nedge](p1)--(p2); \draw[Medge](p2)--(p3); \draw[Nedge](p3)--(p4);
  \node[font=\footnotesize,align=center] at (1.8,-0.8) {starts and ends with $N$:\\ \textbf{augmenting for $M$}};
  \node[v](c1) at (5.4,-0.3){}; \node[v](c2) at (6.6,-0.3){}; \node[v](c3) at (6.6,0.9){}; \node[v](c4) at (5.4,0.9){};
  \draw[Medge](c1)--(c2); \draw[Nedge](c2)--(c3); \draw[Medge](c3)--(c4); \draw[Nedge](c4)--(c1);
  \node[font=\footnotesize,align=center] at (6,-1.1) {cycle: tie};
  \node[v](q1) at (8.2,0){}; \node[v](q2) at (9.4,0){}; \node[v](q3) at (10.6,0){}; \node[v](q4) at (11.8,0){};
  \draw[Medge](q1)--(q2); \draw[Nedge](q2)--(q3); \draw[Medge](q3)--(q4);
  \node[font=\footnotesize,align=center] at (10,-0.8) {starts and ends with $M$:\\ $M$ is ahead here};
  \node[font=\footnotesize,anchor=west] at (0,1.5) {\textcolor{hkblue}{\textbf{---}} edges of $M$ \quad \textcolor{hkorange}{\textbf{- - -}} edges of $N$};
\end{tikzpicture}
\end{center}
{\small Edges in only one of them form \textbf{paths and cycles}. $N$ is bigger overall, so some path has more $N$-edges than $M$-edges: it is augmenting for $M$.}
\end{frame}

% ======================================================================
\begin{frame}{So the classical algorithm is simple}
\begin{columns}[c]
\column{0.5\textwidth}
\begin{algorithm}[H]
\Large
$M\gets\emptyset$\;
\While{an augmenting path exists}{
  find one, flip it
}
\Return{$M$}
\end{algorithm}
\column{0.48\textwidth}
\begin{itemize}\setlength\itemsep{8pt}
  \item Each flip: $+1$ edge. At most $n/2$ flips.
  \item Finding a path: one graph search, $O(m)$.
\end{itemize}
\vspace{2mm}
\begin{alertblock}{Total: $O(mn)$}
about $\frac n2$ rounds, each costing $m$.
\end{alertblock}
\end{columns}
\end{frame}

% ======================================================================
\begin{frame}{But look at what we are wasting}
\begin{columns}[c]
\column{0.45\textwidth}
\centering
\begin{tikzpicture}[scale=1.05]
  \exnodes \exedges
  \draw[aug] (xa)--(y1); \draw[aug] (xb)--(y2); \draw[aug] (xc)--(y4);
  \foreach \v in {xa,xb,xc,xd,xe,y1,y2,y3,y4,y5}{\freering{\v}}
\end{tikzpicture}
\column{0.53\textwidth}
At the start, \emph{every} edge is an augmenting path of length 1.\\[3mm]
Here $a1,\,b2,\,c4$ don't touch, so we could take \textbf{all three at once}.\\[3mm]
Yet the classical algorithm does one full search \emph{per} path.
\begin{alertblock}{Idea}
In each search, grab as many \textbf{non-overlapping} augmenting paths as we can, and flip them together.
\end{alertblock}
\end{columns}
\end{frame}

% ======================================================================
\renewcommand{\curspk}{B}
\begin{frame}{Catch: paths can collide}
\begin{columns}[c]
\column{0.45\textwidth}
\centering
\begin{tikzpicture}[scale=1.05]
  \exnodes \exedges
  \draw[mat] (xa)--(y1); \draw[mat] (xb)--(y2); \draw[mat] (xc)--(y4);
  \draw[aug] (xd)--(y1); \draw[aug] (xa)--(y3);
  \draw[Nedge] (xe)--(y1);
  \freering{xd}\freering{xe}\freering{y3}\freering{y5}
\end{tikzpicture}
\column{0.53\textwidth}
After round 1 the free workers are $d$ and $e$.
\begin{itemize}\setlength\itemsep{4pt}
  \item $d\to1\to a\to3$ \quad\textcolor{hkred}{(red)}
  \item $e\to1\to a\to3$ \quad\textcolor{hkorange}{(orange)}
\end{itemize}
Both need job $1$, worker $a$ and job $3$. Only \textbf{one} can win.
\vspace{2mm}
\begin{block}{Rules for a batch}
\begin{enumerate}
\item Paths must be \textbf{vertex-disjoint}.
\item Use only \textbf{shortest} augmenting paths (the reason comes soon).
\end{enumerate}
\end{block}
\end{columns}
\end{frame}

% ======================================================================
\begin{frame}{Step 1: ripples (BFS) from all free workers}
\begin{columns}[c]
\column{0.62\textwidth}
\centering
\begin{tikzpicture}[scale=0.95]
  \layerA
  \only<2->{\draw[unm] (Ld)--(L1); \draw[unm] (Le)--(L1); \draw[unm] (Ld)--(L2);}
  \only<2->{\draw[mat] (L1)--(La); \draw[mat] (L2)--(Lb);}
  \only<3->{\draw[unm] (La)--(L3); \draw[unm] (Lb)--(L5);}
  \freering{Ld}\freering{Le}
  \only<3->{\freering{L3}\freering{L5}}
  \only<4->{\node[font=\small,text=hkred] at (6.6,-1.6) {first free jobs!};}
\end{tikzpicture}
\column{0.36\textwidth}
\begin{enumerate}\setlength\itemsep{7pt}
  \item<1-> Start with \textbf{all} free workers: layer 0.
  \item<2-> Go along \textbf{non-matching} edges to jobs. Then follow each job's \textbf{matching} edge back to a worker.
  \item<3-> Keep rippling until a layer contains a \textbf{free job}, then \textbf{stop}.
\end{enumerate}
\onslide<4->{The path length so far: $3$. That is the shortest augmenting path length this round.}
\end{columns}
\end{frame}

% ======================================================================
\begin{frame}{Step 2: pick disjoint paths (DFS), skip dead ends}
\begin{columns}[c]
\column{0.62\textwidth}
\centering
\begin{tikzpicture}[scale=0.95]
  \layerA
  \draw[unm] (Ld)--(L1); \draw[unm] (Le)--(L1); \draw[unm] (Ld)--(L2);
  \draw[mat] (L1)--(La); \draw[mat] (L2)--(Lb);
  \draw[unm] (La)--(L3); \draw[unm] (Lb)--(L5);
  \freering{Ld}\freering{Le}\freering{L3}\freering{L5}
  \only<2>{\draw[aug] (Ld)--(L1);}
  \only<3>{\draw[aug] (Ld)--(L1); \draw[aug,line width=3.2pt] (L1)--(La);}
  \only<4->{\draw[aug] (Ld)--(L1); \draw[aug,line width=3.2pt] (L1)--(La); \draw[aug] (La)--(L3);}
  \only<5->{\draw[hkred,very thick] ($(Le)+(-0.3,-0.3)$)--($(Le)+(0.3,0.3)$) ($(Le)+(-0.3,0.3)$)--($(Le)+(0.3,-0.3)$);
            \node[font=\small,text=hkred,align=center] at (3.3,-1.8) {$e$: its only neighbour (job 1) is used.\\ Dead end, never look at it again.};}
\end{tikzpicture}
\column{0.36\textwidth}
\begin{enumerate}\setlength\itemsep{6pt}
  \item<1-> From a free worker, only step \textbf{one layer further}.
  \item<4-> Reach a free job: \textbf{flip} the path. Its vertices are now used.
  \item<5-> Another worker finds everything blocked: mark it \textbf{dead}.
\end{enumerate}
\vspace{2mm}
\onslide<6->{\begin{block}{Why this is fast}
Every edge is tried \textbf{at most once}: one round costs $O(m+n)$.
\end{block}}
\end{columns}
\end{frame}

% ======================================================================
\begin{frame}{Putting it together: the whole run on our example}
\begin{center}
\begin{tikzpicture}[scale=0.55]
  \begin{scope}[shift={(0,0)}]
    \exnodes \exedges
    \draw[aug] (xa)--(y1); \draw[aug] (xb)--(y2); \draw[aug] (xc)--(y4);
    \node[font=\small\bfseries] at (1.4,-3.4) {Round 1};
  \end{scope}
  \begin{scope}[shift={(5.6,0)}]
    \exnodes \exedges
    \draw[mat] (xb)--(y2); \draw[mat] (xc)--(y4);
    \draw[aug] (xd)--(y1); \draw[aug] (xa)--(y1); \draw[aug] (xa)--(y3);
    \node[font=\small\bfseries] at (1.4,-3.4) {Round 2};
  \end{scope}
  \begin{scope}[shift={(11.2,0)}]
    \exnodes \exedges
    \draw[mat] (xa)--(y3); \draw[mat] (xc)--(y4);
    \draw[aug] (xe)--(y1); \draw[aug] (xd)--(y1); \draw[aug] (xd)--(y2);
    \draw[aug] (xb)--(y2); \draw[aug] (xb)--(y5);
    \node[font=\small\bfseries] at (1.4,-3.4) {Round 3};
  \end{scope}
\end{tikzpicture}\\[1mm]
\small\begin{tabular}{lccc}
\toprule
 & \textbf{Round 1} & \textbf{Round 2} & \textbf{Round 3}\\
\midrule
Shortest path length & 1 & 3 & 5\\
Paths flipped & 3 & 1 & 1\\
Matching size & $0\to3$ & $3\to4$ & $4\to5$\\
\bottomrule
\end{tabular}
\end{center}
\alert{Notice:} the shortest path length goes up every round: $1,\,3,\,5$.
\end{frame}

% ======================================================================
\begin{frame}{Why shortest paths? They get used up}
\begin{columns}[T]
\column{0.52\textwidth}
\begin{block}{Fact 1 (Hopcroft--Karp)}
Flipping a shortest augmenting path never creates a \emph{shorter} one. The shortest length only goes \textbf{up}.
\end{block}
\begin{block}{Fact 2}
Two shortest paths of the \emph{same} length never overlap. So every one of them was valid at the start of the round, and we can take them in any order.
\end{block}
\column{0.44\textwidth}
\textbf{Consequence.} After a round that took a \emph{maximal} set of length-$\ell$ paths, none of length $\ell$ is left: if one were, the search would have taken it.
\vspace{3mm}
\begin{alertblock}{So}
After every round, the shortest augmenting path is \textbf{strictly longer}.
\end{alertblock}
\end{columns}
\vspace{3mm}
{\small (Proof sketch in the backup slide. The example shows it: $1\to3\to5$.)}
\end{frame}

% ======================================================================
\begin{frame}{How many rounds? Only about $\sqrt n$}
\centering
\begin{tikzpicture}[x=1.0cm,y=1cm]
  \fill[hkblue!15] (0,0) rectangle (5,1.2);
  \fill[hkorange!22] (5,0) rectangle (10,1.2);
  \node[font=\small] at (2.5,0.6) {first $\sqrt n$ rounds};
  \node[font=\small] at (7.5,0.6) {then at most $\sqrt n$ more};
  \draw[decorate,decoration={brace,amplitude=5pt}] (0,1.35)--(10,1.35) node[midway,above=5pt,font=\small\bfseries] {at most $2\sqrt n$ rounds in total};
\end{tikzpicture}
\vspace{2mm}

\raggedright
\textbf{The argument, in plain words.}
\begin{enumerate}\setlength\itemsep{5pt}
  \item The path length goes up every round. After $\sqrt n$ rounds, every remaining augmenting path has length $\ge\sqrt n$.
  \item Disjoint paths of length $\ge\sqrt n$ use $\ge\sqrt n$ vertices each, and there are only $n$ vertices.
  \item So at most $n/\sqrt n=\sqrt n$ augmentations remain (the best matching differs from ours by that many disjoint paths, as in Berge's picture), and \textbf{each round makes at least one}.
\end{enumerate}
\end{frame}

% ======================================================================
\begin{frame}{The bill}
\begin{block}{Running time}
\[
  \underbrace{O(\sqrt n)}_{\text{rounds}}\ \times\ \underbrace{O(m+n)}_{\text{BFS + DFS per round}}\ =\ O\big((m+n)\sqrt n\big)
\]
\end{block}
\vspace{2mm}
\centering
\begin{tabular}{lrr}
\toprule
 & \textbf{One path per search} & \textbf{Hopcroft--Karp}\\
 & $O(mn)$ & $O(m\sqrt n)$\\
\midrule
$n=10^4,\ m=10^5$ & $\approx10^9$ & $\approx10^7$\\
$n=10^6,\ m=10^7$ & $\approx10^{13}$ & $\approx10^{10}$\\
\bottomrule
\end{tabular}
\vspace{3mm}

{\small Dense graphs ($m\approx n^2$): $O(n^{5/2})$, the bound in the paper's title.}
\end{frame}

% ======================================================================
\begin{frame}{The whole algorithm}
\begin{algorithm}[H]
\Large
$M\gets\emptyset$\;
\Repeat{no free job is reached}{
  \textbf{BFS} from all free workers: build the layers\;
  \textbf{DFS} from each free worker through the layers:\\
  \quad flip every path found, skip dead ends\;
}
\Return{$M$}
\end{algorithm}
\vspace{2mm}
\begin{columns}[T]
\column{0.32\textwidth}\centering\textbf{Berge}\\ \small no augmenting path $\Rightarrow$ done
\column{0.32\textwidth}\centering\textbf{Shortest + disjoint}\\ \small many flips per round
\column{0.32\textwidth}\centering\textbf{$\sqrt n$ rounds}\\ \small lengths grow, few left
\end{columns}
\end{frame}

% ======================================================================
\begin{frame}{What came after}
\begin{itemize}\setlength\itemsep{9pt}
  \item \textbf{Same idea for flows.} A matching is a flow with unit capacities; the rounds are Dinic's \emph{blocking flows}. Even and Tarjan (1975) carried the $O(m\sqrt n)$ bound to unit-capacity networks.
  \item \textbf{Beyond bipartite.} Micali and Vazirani (1980): $O(m\sqrt n)$ for general graphs.
  \item \textbf{Faster in theory.} Recent max-flow work (2022) gives almost-linear time, with much heavier machinery.
  \item \textbf{In practice,} Hopcroft--Karp is still the standard bipartite matching routine: simple and fast.
\end{itemize}
\end{frame}

% ======================================================================
\begin{frame}{Summary}
\begin{columns}[T]
\column{0.55\textwidth}
\begin{enumerate}\setlength\itemsep{8pt}
  \item A bad matching is repaired by a \textbf{chain of swaps}: an augmenting path.
  \item \textbf{No augmenting path $\Leftrightarrow$ maximum} (Berge).
  \item One path per search costs $O(mn)$.
  \item Take \textbf{many shortest, disjoint paths} per search: BFS gives layers, DFS picks paths.
  \item Path lengths keep growing $\Rightarrow$ $O(\sqrt n)$ rounds.
\end{enumerate}
\column{0.42\textwidth}
\begin{alertblock}{$O\big((m+n)\sqrt n\big)$}
the best known in 1973, and still the standard.
\end{alertblock}
\vspace{6mm}
\centering{\Huge\color{hkblue}\textbf{Questions?}}
\end{columns}
\end{frame}

% ======================================================================
\appendix
\begin{frame}{Backup: why shortest paths never get shorter}
\textbf{Claim.} If $P$ is a shortest augmenting path for $M$ and $P'$ is augmenting for $M\oplus P$, then $|P'|\ge|P|+|P\cap P'|$.
\vspace{2mm}

\textbf{Sketch.}
\begin{enumerate}\setlength\itemsep{3pt}
  \item $N=M\oplus P\oplus P'$ is a matching with $|N|=|M|+2$.
  \item By Berge's picture, $M\oplus N=P\oplus P'$ contains \textbf{two disjoint} augmenting paths for $M$, each of length $\ge|P|$ since $P$ is shortest.
  \item So $|P\oplus P'|\ge2|P|$. But $|P\oplus P'|=|P|+|P'|-2|P\cap P'|$. Rearrange. $\square$
\end{enumerate}
\vspace{2mm}
\textbf{Equal length $\Rightarrow$ disjoint.} If $|P'|=|P|$ they share no edge. A shared \emph{vertex} would force a shared \emph{matching edge} (after flipping, every vertex of $P$ is matched along $P$), so they are vertex-disjoint.
\vspace{2mm}

\textbf{Exact count (paper, Thm.\,3).} At most $2\lfloor\sqrt s\rfloor+2$ rounds, where $s$ is the maximum matching size.
\end{frame}

\end{document}
